\newpage
\section{Information Theory}

Information theory is a branch of applied mathematics and electrical engineering involving the quantification of information. 

\subsection{Surprise}

The \emph{surprise} of an event \(X = x\) is defined as 

\[
    \mathrm{Surprise}(x) = S(x) = \ln \left(\frac{1}{P(X = x)}\right)
\]

\subsection{Entropy}

\emph{Entropy} is defined as the expected surprise of a random variable \(X\) 

\[
    \mathrm{Entropy}(X) = H(X) = -\sum_{x \in X} P(X = x) S(x) = -\sum_{x \in X} P(X = x) \ln(P(X = x))
\]

Or in a simpler way

\[
    \mathrm{Entropy}(X) = \sum p(x) \ln\left(\frac{1}{p(x)}\right)
\]

\subsection{Mutual Information}

\emph{Mutual Information} measures how much information one random variable \(X\) contains about another random variable \(Y\). It is defined as

\[
    I(X; Y) = \sum_{x \in X} \sum_{y \in Y} P(X = x, Y = y) \ln \left[\frac{P(X = x, Y = y)}{P(X = x) P(Y = y)} \right]
\]

If one of the random variables is continuous then we can make bins for different range of values and then proceed 
to work with the bins as the discrete classes.

\subsection{Cross Entropy}

\emph{Cross Entropy} is a measure of the difference between two probability distributions for a given random variable or set of 
events.

\[
    H(P, Q) = - \sum_{x \in X} P(X = x) \ln(Q(X = x))    
\]

where \(P\) is the true distribution and \(Q\) is the estimated distribution.   

When there is no discrepancy between both distributions, the cross entropy is equal to the entropy of \(P\).

It also has the following properties:

\begin{itemize}

    \item \emph{Lower Bound:} 
    \[
        H(P,Q) \ge H(P)
    \]

    \item \emph{Asymmetry:}
    \[
        H(P,Q) \ne H(Q, P) 
    \]

\end{itemize}

\subsection{KL-Divergence}

\emph{Kullback-Leibler Divergence} is a measure of how one probability distribution \(Q\) diverges from a second 
expected probability distribution \(P\).

\[
    D_{KL}(P \| Q) = \sum_{x \in X} P(X = x) \ln \left(\frac{P(X = x)}{Q(X = x)}\right)
\]

It measures surprise from using \(Q\) to approximate \(P\). The reverse of the cross entropy, where we use 
\(P\) to approximate \(Q\).

\subsection{Information Gain}

\emph{Information Gain} measures how much information a feature \(A\) gives us about the class label \(Y\). It is defined as

\[
    IG(S, A) = H(S) - \sum_{v \in \text{values}(A)} \frac{|S_v|}{|S|} H(S_v)
\]

where \(H(S)\) is the entropy of the original set \(S\), \(S_v\) is the subset of \(S\) 
for which feature \(A\) has value \(v\), and \(|S|\) and \(|S_v|\) are the sizes of sets \(S\) and \(S_v\) respectively.

In the context of Trees we can think about information gain as  

\[
    IG = H(\text{parent}) - \sum_i  w_i H(\text{children}_i)
\]

where \(w_i\) is the proportion of samples in child \(i\) relative to the parent node. 

\textbf{Example:}

Suppose we have a dataset \(S\) with 10 samples, where 6 belong to class A and 4 belong to class B. The entropy of the original set \(S\) is calculated as:

\[
    H(S) = -\left(\frac{6}{10} \log_2 \left(\frac{6}{10}\right) + \frac{4}{10} \log_2 \left(\frac{4}{10}\right)\right) \approx 0.970
\]

Now, let's say we split the dataset based on a feature \(A\) that can take two values: \(v_1\) and \(v_2\). After the split, we have:

- Subset \(S_{v_1}\) with 4 samples (3 from class A and 1 from class B)
- Subset \(S_{v_2}\) with 6 samples (3 from class A and 3 from class B)

The entropy for each subset is calculated as:

\[
    H(S_{v_1}) = -\left(\frac{3}{4} \log_2 \left(\frac{3}{4}\right) + \frac{1}{4} \log_2 \left(\frac{1}{4}\right)\right) \approx 0.811
\]

\[
    H(S_{v_2}) = -\left(\frac{3}{6} \log_2 \left(\frac{3}{6}\right) + \frac{3}{6} \log_2 \left(\frac{3}{6}\right)\right) = 1.0    
\]  

The weighted average entropy after the split is:

\[
    \sum_{v \in \text{values}(A)} \frac{|S_v|}{|S|} H(S_v) = \frac{4}{10} \cdot H(S_{v_1}) + \frac{6}{10} \cdot H(S_{v_2}) \approx 0.4 \cdot 0.811 + 0.6 \cdot 1.0 \approx 0.924
\]

Therefore, the information gain from splitting on feature \(A\) is:

\[
    IG(S, A) = H(S) - \sum_{v \in \text{values}(A)} \frac{|S_v|}{|S|} H(S_v) \approx 0.970 - 0.924 = 0.046
\]  

\subsection{Gini Impurity}

\emph{Gini Impurity} is a measure of how often a randomly chosen element from the set would be incorrectly labeled if it was randomly 
labeled according to the distribution of labels in the subset. It is calculated as

\[
    Gini(S) = 1 - \sum_{i=1}^{C} p_i^2
\]

where \(C\) is the number of classes and \(p_i\) is the proportion of instances in class \(i\) in the set \(S\).

The closer the Gini impurity is to 0, the more pure the set is.


